% =====================================================================
%  Conferenze, Zagarolo (On The Road), giorno 1 — Intervista 5
%  Publishing indie, fiducia e ``magia'': il modello Raw Fury
%  Fonte: trascrizione automatica "Track 13". Intervista in inglese a
%  Jonas, cofondatore del publisher svedese Raw Fury, condotta da un
%  organizzatore di eventi per studenti, con domande dal pubblico.
%  Etichette con prefisso "c2-5:". Nessun pacchetto aggiuntivo.
% =====================================================================

\chapter[Publishing indie e originalità: Raw Fury]{Publishing indie, fiducia e
``magia'': l'esperienza di Raw Fury}
\label{cap:c2-5-raw-fury}

Dopo il marketing (capitolo~\ref{cap:c2-1-marketing-publisher}) e la produzione
(capitolo~\ref{cap:c2-2-producer}) visti dall'interno di un publisher, questa
intervista guarda al publishing dal punto di vista di chi un publisher l'ha
\emph{fondato}. L'ospite racconta perché ha creato un'etichetta costruita sulla
fiducia invece che sulla trattativa, come sceglie i giochi da firmare, che cosa
pensa dell'intelligenza artificiale applicata allo scouting e che cosa si aspetta
da un pitch. Il capitolo è organizzato così:
\begin{enumerate}
  \item l'ospite e il suo percorso (sezione~\ref{sec:c2-5-ospite});
  \item come è cambiata l'industria e perché servono i publisher
        (sezione~\ref{sec:c2-5-industria});
  \item un publisher fondato sulla fiducia: il contratto unico e il
        \emph{flywheel} (sezione~\ref{sec:c2-5-fiducia});
  \item l'IA come livello di astrazione e il rischio dello scouting automatico
        (sezione~\ref{sec:c2-5-ia});
  \item istinto ed esperienza: come si riconosce la ``magia''
        (sezione~\ref{sec:c2-5-magia});
  \item le domande del pubblico: pitch, gusto, tecnica, wishlist, community
        (sezione~\ref{sec:c2-5-domande}).
\end{enumerate}

% ---------------------------------------------------------------------
\section{L'ospite}
\label{sec:c2-5-ospite}

L'ospite è \textbf{Jonas}\footnote{La trascrizione automatica riporta il cognome
in forme diverse e incompatibili (``Jonas Atkinson'', ``Jonas Santos''). Dal
contesto (islandese, cofondatore di Raw Fury) si tratta con ogni probabilità di
\textbf{Jónas Antonsson}; il cognome va comunque considerato una ricostruzione.},
islandese, cofondatore del publisher svedese \textbf{Raw
Fury}\footnote{Il nome dell'azienda è storpiato in molti modi nella trascrizione:
``Brocury'', ``Row Fury'', ``Rocky Rain'', ``Rocky League'', ``Rothfury''. Il nome
corretto è \emph{Raw Fury}.}. Oggi siede nel \textbf{consiglio di
amministrazione} e non ha più ruoli operativi nell'azienda.

Il suo percorso:
\begin{itemize}
  \item dai primi anni Duemila ha guidato per circa \textbf{dieci anni} uno studio
        di sviluppo in Islanda;
  \item si è poi trasferito in \textbf{Svezia}, dove ha fondato un publisher;
  \item conta circa \textbf{venticinque anni} di esperienza nell'industria, che
        nel corso dell'intervista indica come la base del suo ``istinto''
        (sezione~\ref{sec:c2-5-magia}).
\end{itemize}

% ---------------------------------------------------------------------
\section{Come è cambiata l'industria}
\label{sec:c2-5-industria}

\subsection{Dalle sfide in serie alle sfide in parallelo}

Secondo Jonas, un tempo le grandi trasformazioni del settore arrivavano
\textbf{una alla volta}: c'era il tempo di assorbirne una prima che arrivasse la
successiva. Oggi arrivano \textbf{in parallelo}. Cita in particolare:
\begin{itemize}
  \item il passaggio alla \textbf{distribuzione digitale};
  \item l'esplosione del \textbf{mobile};
  \item la \textbf{democratizzazione dei motori di gioco}, che ha reso possibile a
        chiunque sviluppare un gioco.
\end{itemize}

\subsection{Dall'oceano blu all'oceano rosso}

Intorno al \textbf{2015} il mercato mobile era passato da \textbf{oceano blu}
(spazio poco affollato, con molte opportunità) a \textbf{oceano rosso} (mercato
saturo, dove si compete al ribasso). Jonas era convinto che la stessa cosa sarebbe
successa \textbf{su tutte le piattaforme}: quando chiunque può pubblicare, i giochi
diventano troppi per essere notati da soli. Per questo, secondo lui, ci sarebbe
stato \textbf{bisogno di publisher}.

\subsection{Self-release non è self-publishing}

In quegli anni circolava il mantra ``\textbf{tutti possono autopubblicarsi}''.
Per Jonas è una \textbf{bugia}: tutti possono fare \emph{self-release}, cioè
\textbf{mettere in vendita} il proprio gioco su un negozio digitale; ma
\emph{pubblicare} significa molto di più (marketing, posizionamento, rapporti con
le piattaforme, localizzazione, porting, finanziamento), e questo non lo può fare
chiunque.

% ---------------------------------------------------------------------
\section{Un publisher fondato sulla fiducia}
\label{sec:c2-5-fiducia}

\subsection{Il problema dei publisher tradizionali}

I publisher dell'epoca erano \textbf{tarati sui grandi team}: strutture, processi e
contratti (Jonas ricorda contratti da \textbf{52 pagine}) pensati per studi con
decine di persone e budget importanti. Applicati a un team indipendente di poche
persone, diventavano sproporzionati.

L'idea di Jonas era un publisher fondato sulla \textbf{relazione} e sulla
\textbf{fiducia}. La sua lettura dei piccoli team è netta:
\begin{quote}
i team piccoli sono pessimi nel fare budget e nel gestire lo \emph{scope}, ma sono
brillanti nel \emph{fare}.
\end{quote}
Il compito del publisher, quindi, è coprire le debolezze (budget, pianificazione,
tutto ciò che sta intorno al gioco) lasciando al team ciò che sa fare meglio.

\subsection{Il primo contratto e il flywheel}

Il primo contratto fu scritto insieme a \textbf{Thomas}\footnote{Thomas van den
Berg, creatore di \emph{Kingdom}; il cognome non è pronunciato chiaramente nella
trascrizione.}, il creatore di \textbf{\emph{Kingdom}}. Il modello economico è
quello del \textbf{flywheel} (volano): il \textbf{catalogo} dei giochi già usciti
genera i ricavi che \textbf{finanziano i giochi successivi}, e ogni nuovo successo
alimenta ulteriormente il ciclo.

\subsection{Gli stessi termini per tutti}

Raw Fury offre gli \textbf{stessi termini contrattuali a tutti gli sviluppatori}.
Le conseguenze, nelle parole di Jonas:
\begin{itemize}
  \item la relazione parte dalla \textbf{fiducia}, senza una trattativa ``a somma
        zero'' in cui ciò che guadagna una parte lo perde l'altra;
  \item gli sviluppatori possono \textbf{parlarsi tra loro} apertamente, perché
        nessuno ha condizioni segrete da proteggere;
  \item nasce una vera \textbf{community} di sviluppatori intorno all'etichetta.
\end{itemize}
Il contratto è stato poi \textbf{reso pubblico}. Un avvocato lo ha analizzato in un
video su YouTube, e oggi molti publisher indipendenti lo usano come
\textbf{modello} per i propri.

% ---------------------------------------------------------------------
\section{L'IA come livello di astrazione}
\label{sec:c2-5-ia}

\subsection{Un'altra astrazione nella storia dell'informatica}

Jonas legge l'intelligenza artificiale come l'ennesimo \textbf{livello di
astrazione} nella storia della programmazione: dall'\textbf{assembly} a
\textbf{Fortran} e \textbf{COBOL}, poi \textbf{C}, poi \textbf{Java} e
\textbf{C++}. Ogni passaggio ha allontanato il programmatore dalla macchina e
cambiato il lavoro, senza eliminarlo.

\subsection{Il rischio: lo scouting automatico}

Il pericolo che vede è un altro: che i publisher usino l'IA per \textbf{filtrare i
pitch}. In quel caso:
\begin{itemize}
  \item tutti i publisher finirebbero per vedere \textbf{gli stessi 10--20 giochi},
        quelli che il modello giudica promettenti;
  \item l'IA \textbf{predice a partire dal passato}: premia ciò che somiglia a ciò
        che ha già funzionato;
  \item di conseguenza \textbf{non avrebbe mai firmato} un gioco come
        \textbf{\emph{Blue Prince}}\footnote{Nella trascrizione ``Google
        Blueprints'': si tratta di \emph{Blue Prince}, il puzzle-mystery con
        elementi roguelike pubblicato da Raw Fury nel 2025.}, che non somigliava
        a nulla.
\end{itemize}

\subsection{Le persone sono pigre}

Il problema, dice, è che le persone sono \textbf{pigre}, come i giocatori che, se
esiste una scorciatoia, la prendono anche quando rovina il divertimento. Cita uno
studio sull'apprendimento della \textbf{matematica} prima e dopo l'introduzione
dell'IA, da cui trae la formula:
\begin{quote}
l'IA è il miglior insegnante e il peggior tutor.
\end{quote}
Usata per capire, insegna moltissimo; usata per farsi fare i compiti, impedisce di
imparare.

% ---------------------------------------------------------------------
\section{Istinto ed esperienza: riconoscere la ``magia''}
\label{sec:c2-5-magia}

\subsection{Il capitale e i fogli Excel}

``Il capitale odia ciò che non entra in un foglio Excel'', dice Jonas. Ma quella
dei videogiochi è un'industria \textbf{guidata dalle persone}, in cui le scelte più
importanti non si riducono a una tabella.

Come scout, Jonas ha avuto un \textbf{hit rate del 76\%}, inteso come percentuale
di giochi firmati che sono almeno \textbf{andati in pareggio}. Per confronto, cita
il \textbf{57\%} del miglior giocatore di poker. Il suo criterio è semplice:
\begin{quote}
quando sento la magia.
\end{quote}
È \textbf{istinto}, ma un istinto costruito su \textbf{venticinque anni di
esperienza}.

\subsection{Tre segnali}

Alla richiesta di rendere più concreto quel ``sentire la magia'', indica tre
segnali.
\begin{enumerate}
  \item \textbf{Innamorarsi di una meccanica che interessa solo a te.} Un
        creatore ossessionato da un'idea che a nessun altro sembra importante è
        spesso all'origine dei giochi più originali.
  \item \textbf{Le persone che si sporgono in avanti.} Osservando chi prova il
        gioco, conta il gesto di \emph{lean in}, il coinvolgimento fisico. Alle
        fiere, se \textbf{più del 5\%} dei giocatori \textbf{torna allo stand}
        dopo mezz'ora o il giorno dopo, è un segnale fortissimo.
  \item \textbf{I pattern nei pitch.} Ricevendo migliaia di pitch, si vedono le
        tendenze prima che arrivino sul mercato:
        \begin{itemize}
          \item prima del \textbf{Covid} i pitch erano sempre più
                \textbf{cupi e nichilisti};
          \item circa \textbf{sei mesi dopo} sono arrivati palette pastello,
                giardini, temi rassicuranti: la nascita dei \textbf{cozy game};
          \item dopo il successo di \textbf{\emph{Vampire Survivors}}, \textbf{un
                pitch su tre} era ``\emph{Vampire Survivors} più qualcosa''.
        \end{itemize}
\end{enumerate}

\subsection{Team esperto o gioco straordinario?}

Meglio un team esperto con un gioco normale o un team giovane con un gioco
straordinario? Jonas sceglie il \textbf{team giovane}, se nel gioco c'è la magia.
Un team inesperto sbaglierà budget e tempi, ma è proprio questo che il publisher
deve saper gestire: il talento e l'idea non si possono aggiungere dall'esterno.

% ---------------------------------------------------------------------
\section{Le domande del pubblico}
\label{sec:c2-5-domande}

\subsection{Basta una frase per presentare un gioco?}

No. Una frase sola non basta, anche se aiuta: \emph{Blue Prince}, per esempio, si
riassumeva in ``un mystery con meccaniche roguelike''. L'\textbf{ordine} in cui
Jonas guarda il materiale di un pitch è:
\begin{enumerate}
  \item un \textbf{video di gameplay} di circa \textbf{due minuti}, meglio se con
        lo sviluppatore che \textbf{gioca e commenta};
  \item una \textbf{build giocabile};
  \item solo dopo, il \textbf{pitch deck} e l'email.
\end{enumerate}
Non ha \textbf{mai scartato un gioco per il deck}: un deck brutto non affossa un
buon gioco, mentre un buon video può convincere da solo.

\subsection{Esiste un gusto oggettivo?}

No. Jonas porta un esempio volutamente bizzarro: esistono appassionati della
\textbf{moquette degli alberghi}, che la fotografano e la collezionano. Qualunque
cosa può avere il suo pubblico.

Tuttavia alcuni gusti coprono una \textbf{porzione più ampia del diagramma di
Venn} dei gusti delle persone. Cita \textbf{MySpace} e la cosiddetta
\textbf{``Theory of Ugly''}: pagine personali bruttissime ma
\textbf{autentiche}, che attiravano proprio per questo. L'\textbf{autenticità
attrae}.

\subsection{Buon gusto ma tecnica debole}

Se un team ha una visione forte ma competenze tecniche limitate, i publisher
\textbf{full-stack} possono aiutare: per esempio a \textbf{scalare un
multiplayer}. Raw Fury ha firmato un gioco bellissimo sapendo fin dall'inizio che
la parte più difficile sarebbe stata \textbf{ottimizzarlo per le console}: lo ha
fatto perché quel lavoro poteva coprirlo il publisher, mentre la bellezza del gioco
no.

\subsection{Bisogna dimostrare la domanda con le wishlist?}

Molti publisher chiedono di presentarsi con numeri di \textbf{wishlist} già alti.
Jonas vede il rischio dell'\textbf{uovo e della gallina}: se hai già trazione,
perché dovresti cercare un publisher? Fa il paragone con i \textbf{venture
capitalist}, che chiedono trazione prima di investire. Per lui
\textbf{costruire la trazione è compito del publisher}.

Richiama il \textbf{Principio di Dilbert} di \textbf{Scott Adams}, cioè il lavoro
che sembra importante ma non produce nulla, e ricorda che gli \textbf{scout
musicali} nei club (in Germania) hanno trovato band come i \textbf{Beatles} e i
\textbf{Guns N' Roses} ascoltandole dal vivo, non guardando i loro numeri.
Spostare tutta la responsabilità sugli sviluppatori, conclude, porta a
\textbf{perdere arte e creatività}.

\subsection{La community}

Ogni sviluppatore dovrebbe costruire la propria \textbf{community}. La domanda da
farsi è brutale: ``\emph{Why should anyone fucking care?}'', perché a qualcuno
dovrebbe importare del tuo gioco? Il consiglio:
\begin{itemize}
  \item mettere il gioco \textbf{in mano alle persone presto};
  \item usare la \textbf{voce autentica del creatore}, che nessun publisher può
        replicare;
  \item ricordare che la community \textbf{appartiene allo sviluppatore}, non al
        publisher.
\end{itemize}

\subsection{Adattare il proprio stile al publisher?}

In parte. Esistono publisher dal profilo \textbf{stretto}, come
\textbf{Devolver}, dove ``Tim''\footnote{Nome riportato così nella
trascrizione, senza cognome; si riferisce a una persona di Devolver che conosce a
fondo il tipo di giochi pubblicati dall'etichetta.} conosce a fondo un genere
preciso di giochi, e publisher dal profilo \textbf{largo}, come Raw Fury,
\textbf{eclettico} ma sempre in cerca di ``magia'' e di un taglio artistico.

Il consiglio pratico: \textbf{niente invii a pioggia}. Meglio scegliere i
\textbf{5--10 publisher} più adatti e preparare per ognuno un \textbf{pitch
personalizzato}. In ogni caso, ciò che conta di più resta il \textbf{video}.

% ---------------------------------------------------------------------
\section*{Punti chiave del capitolo}
\addcontentsline{toc}{section}{Punti chiave del capitolo}

\begin{itemize}
  \item Le trasformazioni dell'industria (digitale, mobile, motori accessibili)
        arrivano \textbf{in parallelo}; con mercati saturi servono i publisher:
        \textbf{self-release non è self-publishing}.
  \item Raw Fury nasce come publisher fondato su \textbf{relazione e fiducia}, con
        \textbf{gli stessi termini per tutti} e un contratto reso pubblico; il
        catalogo finanzia i giochi futuri (\textbf{flywheel}).
  \item L'IA è un nuovo livello di astrazione; usata per filtrare i pitch
        appiattirebbe lo scouting sul passato e scarterebbe i giochi originali.
  \item Scouting: \textbf{istinto più esperienza}; segnali sono l'ossessione per
        una meccanica, il pubblico che ``si sporge'' e torna allo stand, i pattern
        nei pitch.
  \item Pitch: prima un \textbf{video di gameplay} di due minuti, poi la build,
        poi il deck; pitch personalizzati a pochi publisher scelti.
  \item Nessun gusto oggettivo, ma l'\textbf{autenticità attrae}; la community è
        dello sviluppatore, e costruire la trazione è compito del publisher.
\end{itemize}
